% Pregunta 479 (2014-septiembre-OEP2014, n.º 11): equilibrio de Stackelberg con
% demanda lineal. Ejes en unidades de (a-c)/b.
%   R_2: q_2 = 1/2 - q_1/2 (la regla de la opción d, pendiente -1/2)
%   R_1 (Cournot): q_1 = 1/2 - q_2/2, de (1/2, 0) a (0, 1); Cournot C = (1/3, 1/3)
%   Isobeneficio del líder de nivel k: q_1 (1 - q_1 - q_2) = k
%     -> q_2 = 1 - q_1 - k/q_1; con k = 1/8 su pendiente, -1 + 1/(8 q_1^2),
%        vale -1/2 en q_1 = 1/2: tangente a R_2 en S = (1/2, 1/4)
% Autor: Víctor Gutiérrez Marcos
\documentclass[tikz,border=6pt]{standalone}
\input{tikz-tcee}
\begin{document}
\begin{tikzpicture}[tcee,x=5.4cm,y=5.4cm,
    declare function={iso(\k,\q)=1-\q-\k/\q;}]
  \ejes{$q_1$}{$q_2$}{1.22}{1.08}
  % R_1 de Cournot, de referencia
  \draw[thick,dotted,tceegris] (0.5,0) -- (0,1);
  \node[etiqueta,text=tceegris,anchor=west] at (0.03,0.96) {$R_1$ (Cournot)};
  % isobeneficio del líder tangente a R_2 en S
  \draw[curva4] plot[domain=0.1465:0.8535,samples=80] (\x,{iso(1/8,\x)});
  \node[etiqueta,text=tceedorado!80!black,anchor=west,align=left,font=\footnotesize]
    at (0.72,0.30) {isobeneficio\\ del líder};
  % R_2: función de reacción del seguidor
  \draw[curva1] (0,0.5) -- (1,0);
  \node[etiqueta,text=tceemorado,anchor=west,align=left] at (0.66,0.57)
    {$R_2$: $q_2=\frac{a-c}{2b}-\frac12\,q_1$\\[1pt] {\footnotesize (pendiente $-\tfrac12$)}};
  \draw[thin,tceemorado,->] (0.74,0.485) -- (0.80,0.115);
  \node[etiqueta,left] at (0,0.5) {$\frac{a-c}{2b}$};
  \node[etiqueta,below] at (1,0) {$\frac{a-c}{b}$};
  % Cournot
  \node[punto,fill=tceegris,label={[etiqueta,text=tceegris]above right:$C$}] at (1/3,1/3) {};
  % Stackelberg
  \node[punto,fill=tceerojo,inner sep=1.8pt,
        label={[etiqueta,text=tceerojo,fill=white,inner sep=1pt]above right:$S$}] (S) at (0.5,0.25) {};
  \proyeccion{(S)}{$\frac{a-c}{2b}$}{$\frac{a-c}{4b}$}
  \node[etiqueta,text=tceerojo,anchor=north west,align=left,font=\footnotesize]
    at (0.36,0.95) {en $S$, el líder produce el doble\\ que el seguidor y gana el doble:\\[2pt] $\pi_1=\frac{(a-c)^2}{8b}>\pi_2=\frac{(a-c)^2}{16b}$};
\end{tikzpicture}
\end{document}
